\chapter{Glossary, help and reporting problems}
\label{app:glossary}

The tutorial chapters use their words without ceremony and expect you to keep
moving. This appendix is the one to keep beside the machine: what the vocabulary
means, where the program sends you when you ask for help, how to report a fault
so that it can actually be fixed, which hardware is currently known to work,
and a checklist to run through before the beam comes on.

\section*{Glossary}

Terms are listed as an operator meets them, in alphabetical order. Where a term
has a specific meaning inside MeerK40t --- rather than the general laser-shop
meaning --- that meaning is the one given.

{\small
\begin{longtable}{@{}p{3.3cm}p{11.1cm}@{}}
\toprule
\textbf{Term} & \textbf{Meaning} \\
\midrule
\endfirsthead
\toprule
\textbf{Term} & \textbf{Meaning} \\
\midrule
\endhead
\bottomrule
\endlastfoot

\textbf{Bed} & The flat working area of the machine, and the white rectangle
MeerK40t draws for it in the scene. Its size comes from the device's
\setting{bedwidth} and \setting{bedheight} settings, so a bed with the wrong
proportions is a device-configuration problem, not a drawing problem
(Chapter~\ref{ch:connect}). \\

\textbf{Blob} & The device-specific data that a plan is converted into before it
is streamed to the controller --- the \cmd{blob} stage of planning. A blob shows
up as a \ui{blob} node in the \panel{Tree}, and can be executed, converted back
to elements, or deleted. \\

\textbf{Burn} & Running a job with the beam on: the machine is marking or cutting
material. Inside MeerK40t ``burn'' also names the work itself --- the
\tabname{Burn-Operation} tab of the \panel{Tree} lists what will burn and in what
order, and the status bar reports burn progress. \\

\textbf{Classification} & The rule that decides which operation an element
belongs to, by comparing the element's stroke and fill colours with the colours
of your operations. Whether matching is exact or approximate, whether unmatched
elements go to a ``default'' operation, and whether MeerK40t may create a new
operation for them are all set on the \tabname{Classification} tab of the
Preferences window (Chapter~\ref{ch:ops}). \\

\textbf{Cutcode} & MeerK40t's internal list of laser primitives --- lines, arcs,
curves, raster sweeps, dwells and moves --- each carrying the settings it must be
burned with. Operations are turned into cutcode, and the drivers consume cutcode;
you never edit it directly. \\

\textbf{DPI} & Dots per inch: how many raster scan lines fit into an inch of
travel. It is the reciprocal of the raster step, so doubling the DPI halves the
step, roughly doubles the burn time, and only pays off while the step is still
larger than your laser's spot (Chapter~\ref{ch:tuning}). \\

\textbf{Dwell} & A deliberate pause at one point with the beam on, used to mark a
dot or to pierce material. MeerK40t's \mkseries{Dots} operation carries a dwell
time and produces one dwell per point element attached to it. \\

\textbf{Effect} & A node that rewrites geometry instead of burning something
itself. \emph{Hatch} fills a closed shape with parallel lines, \emph{Wobble}
follows the outline with a wavy contour, and \emph{Warp} distorts geometry into a
grid. Effects can be attached to elements or appended to a \mkseries{Cut} or
\mkseries{Engrave} operation. \\

\textbf{Element} & One piece of artwork in the \panel{Tree}: a path, rectangle,
ellipse, line, polyline, point, text, image or blob. Elements live in the
\ui{Elements} branch; operations do not contain them, they hold references to
them. \\

\textbf{Emphasis (selected)} & The emphasised nodes are the current selection ---
the ones the toolbar, the property panes and the console act on. Emphasising an
element also highlights the operations that reference it, and emphasising an
operation highlights its sub-operations and the elements it uses. \\

\textbf{Engrave} & A vector operation that traces a path with the beam on but
does not cut through: thin, dark lines, good for text and detail. The same
geometry burned with a \mkseries{Cut} operation and heavier settings is a cut
(Chapter~\ref{ch:ops}). \\

\textbf{Frame / outline} & Running the machine around the outside of a job so you
can see where it will land before burning it. In MeerK40t the
\cmd{outline <offset>} command --- also \menu{Outline element(s)...} on an element
--- creates an offset outline path around the emphasised elements, and some
drivers add machine-side frame buttons (Newly: \btn{Draw Frame},
\btn{Move Frame}). \\

\textbf{Galvo} & A marking head that steers the beam with mirrors across a fixed
field instead of moving a gantry, as on fibre and UV marking machines. MeerK40t
drives Ezcad2-compatible JCZ galvo controllers through the Balor driver. \\

\textbf{G-code} & The motion language used by GRBL-family controllers. MeerK40t
generates G-code for GRBL devices and can export it as \file{.gcode},
\file{.nc} or \file{.gc}; K40, Ruida, Moshi and Newly machines speak their own
protocols instead. \\

\textbf{Hatch} & A fill made of parallel lines inside a closed shape, with a
hatch distance and angle (algorithms: scanline, Eulerian, spiral). Hatching is
how an area is engraved in sweeps instead of outline by outline. \\

\textbf{Home} & A fixed position the machine can find for itself, usually with
limit switches; the \cmd{home} command and the controller's own Home button send
the head there. Home is the machine's reference point and is not necessarily the
origin a job is placed from (Chapter~\ref{ch:connect}). \\

\textbf{Image operation} & The \mkseries{Image} operation, which engraves an
existing image pixel by pixel with the operation's own settings. It is not the
same as \mkseries{Raster}, which first renders everything attached to it into one
intermediate image and then burns that. \\

\textbf{Job} & One compiled unit of work: the operations built into cutcode,
wrapped up and handed to the spooler to stream to the machine. Jobs appear in the
\panel{Spooler} pane with their status, passes, runtime and estimate. \\

\textbf{Job origin (job start, user origin)} & The point a job is placed from, so
that the same artwork can be burned at several positions on the bed. In MeerK40t
it is a \emph{placement} node under the operations branch
(\menu{Define Job Start Position}), and you may define several. \\

\textbf{Kerf} & The width of material that the beam destroys. A part drawn
20\unit{mm} wide comes out narrower unless you compensate, which is why operations
expose a kerf value and MeerK40t can offset paths by half of it. \\

\textbf{Laser path} & The preview of where the beam will actually travel, drawn
over the artwork with the direction of each cut. Toggle it under
\menu{View>Scene Appearance}. \\

\textbf{Lid / interlock} & The safety switch that stops the beam when the
enclosure is opened. MeerK40t cannot see it and cannot override it; a machine
whose lid switch has been defeated is a machine that will injure someone. \\

\textbf{Lihuiyu} & The driver family for Lihuiyu M2-Nano and M3-Nano boards ---
the original K40 controller. USB only, vector and raster both supported. See
\ref{app:drivers} for its settings. \\

\textbf{Machine origin} & The zero point the machine itself works in, fixed by
homing or by the controller. Where it physically sits --- front-left, centre, or
somewhere else --- is decided by your device profile and controller, so establish
it with a light test on scrap before you cut anything valuable. \\

\textbf{Moshi} & The driver for Moshiboard controllers, found in older
laser machines. USB, with the board's own command set. \\

\textbf{Newly} & The driver for NewlyDraw System 8.1 machines. USB, with its own
file slots on the controller and \btn{Draw Frame} / \btn{Move Frame} controls for
checking placement. \\

\textbf{Op} & A single operation node in the \panel{Tree}, holding one recipe:
its type, colour, speed, power, passes and effects. An op is what you select,
edit and drag elements onto. \\

\textbf{Operation} & The recipe for burning a group of artwork: cut, engrave,
raster, image or dots, plus its settings. Operations are listed in the
\ui{Operations} branch and are planned in order, top to bottom
(Chapter~\ref{ch:ops}). \\

\textbf{Overscan} & Extra travel added at each end of a raster line so that the
head is already at full speed when it reaches the image. Without it, the ends of
each sweep linger and burn darker. \\

\textbf{Passes} & How many times an operation repeats the same path or sweep.
Two passes at moderate power usually cut cleaner than one at maximum, especially
in plywood and acrylic. \\

\textbf{PPI} & Pulses per inch: how many laser pulses are emitted per inch of
travel, shown as \emph{frequency} in kilohertz on machines that expose it.
It changes how an edge looks more than whether the material cuts at all. \\

\textbf{Plan} & The working copy of the selected operations that the planner
builds, stage by stage, and then hands to the spooler. Plans are numbered, so you
can hold several at once; the console's \cmd{plan} commands act on them
(Chapter~\ref{ch:order}). \\

\textbf{Preopt / optimize} & The two optimisation stages of planning.
\cmd{preopt} prepares the plan for reordering and adds the optimiser passes;
\cmd{optimize} runs them, reducing travel and cutting inner contours before outer
ones. With optimisation switched off, a plan skips both. \\

\textbf{Profile} & A saved machine configuration --- a device entry with its
driver, bed size, origin and settings. Profiles are chosen and edited in the
\panel{Devices} pane; the device chooser at first start offers generic profiles
grouped under each driver. \\

\textbf{Raster} & Pixel artwork, and the operation that burns it line by line. A
\mkseries{Raster} operation renders everything attached to it into one image and
sweeps that image; a raster engrave is much slower than a vector one over the
same area (Chapter~\ref{ch:engrave}). \\

\textbf{Reference} & A node under an operation that points at an element in the
elements branch, which is how one element can belong to several operations
without being copied. (The alignment ``reference object'' --- an element other
elements are aligned to --- is a different thing with the same name.) \\

\textbf{Regmark} & A registration mark: artwork used only to align the work ---
with a camera, a jig, or the machine's own pointer. Regmarks live in their own
\ui{Regmarks} branch so they are never burned, and invisible objects in imported
files can be loaded there. \\

\textbf{Rotary} & A rotating attachment that turns a cylindrical workpiece so the
laser can engrave around it. MeerK40t configures rotary work per device and also
has support for cylindrical material on flat beds. \\

\textbf{Ruida} & The driver for Ruida controllers, the black-and-yellow
CO\textsubscript{2} machines with an RDC6442- or RDC6445-class panel. Network
first, and able to emulate a Ruida controller for other software. \\

\textbf{Spooler} & The queue of jobs waiting for the machine, and the pane that
shows it (\panel{Spooler}, docked at the bottom of the main window on first use).
It reports each job's items, steps, passes, runtime and estimate, and carries the
\btn{Pause} and \btn{Abort} controls. \\

\textbf{Step (raster step)} & The distance between adjacent raster scan lines,
in millimetres or inches. It is the number MeerK40t actually works in; DPI is the
same quantity expressed as lines per inch, and around 500\unit{dpi} the step is
about 0.05\unit{mm}. \\

\textbf{SVG} & The vector file format MeerK40t imports best, and also its native
project format: \menu{File>Save} writes your project --- artwork, operations and
notes --- out as SVG. \\

\textbf{TTF / Hershey} & The font kinds available to vector (line) text:
TrueType (\file{.ttf}), Hershey (\file{.jhf}) and Autocad shape (\file{.shx})
fonts. All three are converted to paths that can be scaled and cut like any other
vector shape (Chapter~\ref{ch:engrave}). \\

\textbf{Vector} & A shape described mathematically --- a line, arc or curve ---
rather than as pixels. Vectors scale without degrading and are burned by tracing
the path, which is why they cut and score. \\

\textbf{Wobble} & An effect that makes the beam follow a wavy contour around a
path --- circle, sine wave, sawtooth, jigsaw, gear and other shapes --- with a
radius, an interval and a speed. It produces decorative, stitch-like edges and
wider, softer cuts. \\

\textbf{Wordlist} & A named list of text values (optionally CSV-backed) whose
entries are substituted into text at burn time, advancing one entry per burn.
Used for serial numbers, batch labels and the like. \\

\textbf{Work area} & The region the machine can actually reach, which is the bed
minus whatever the head, the gantry or a rotary attachment takes away. Elements
outside the work area are either unburnable or will be cut somewhere unexpected;
the \panel{Tree} warns about them. \\

\textbf{Zero, 0,0} & The coordinate origin the scene measures positions from. A
\ui{Goto} utility node placed at 0,0 is shown in the \panel{Tree} as
\ui{Origin}, which is a convenient way to see where the zero point of a job
sits. \\

\end{longtable}
}

\section*{Getting help}

MeerK40t is developed by a community that mostly answers questions in public
forums rather than by email, and the \menu{Help} menu is the front door to those
places. The entries are always the same; only the addresses they open change.

\begin{center}\small
\begin{tabular}{@{}p{5.0cm}p{9.3cm}@{}}
\toprule
\textbf{Menu entry} & \textbf{What it opens} \\
\midrule
\menu{Help>Help} (\key{F1}) and
\menu{Help>Online Reference} (\key{F11}) & Context-sensitive online help.
MeerK40t looks at the widget under the mouse pointer, walks up through its
parents until it finds help text, and opens the wiki page for that section;
if nothing under the pointer has help text, it opens the general interface
section instead. \\
\menu{Help>Beginners' Help} & The wiki index for new users:
{\scriptsize\file{https://github.com/meerk40t/meerk40t/wiki/Beginners:-0.-Index}}. Start here
when the manual has not answered the question. \\
\menu{Help>Github} & The project page:
{\scriptsize\file{https://github.com/meerk40t/meerk40t}}. \\
\menu{Help>Releases} & The release list, where installers, stable releases and
weekly builds live:
{\scriptsize\file{https://github.com/meerk40t/meerk40t/releases}}. \\
\menu{Help>Facebook} & The K40 MeerK40t group:
{\scriptsize\file{https://www.facebook.com/groups/716000085655097}}. \\
\menu{Help>Discord} & The project's Discord server, where developers and users
talk:
{\scriptsize\file{https://discord.gg/vkDD3HdQq6}}. \\
\menu{Help>Makers Forum} & The MeerK40t section of the Makers Forum:
{\scriptsize\file{https://forum.makerforums.info/c/k40/meerk40t/120}}. \\
\menu{Help>Feature request/feedback} & The GitHub Discussions chooser for
reporting requests and ideas:
{\scriptsize\file{https://github.com/meerk40t/meerk40t/discussions/new/choose}}. Feature
requests belong here rather than in the issue tracker. \\
\menu{Help>Check for Updates} & Runs the update check immediately and reports the
result in an \ui{Update-Info} dialog. \\
\menu{Help>Tips \&\& Tricks} & The Tips window: previous and next tip, a
\btn{Try it out} button, and switches for showing tips at startup. \\
\menu{Help>About MeerK40t} & The About window, described below. \\
\bottomrule
\end{tabular}
\end{center}

\begin{notebox}[Other doors into the same project]
On macOS the first two entries differ: \btn{MeerK40t Help} opens a local help
file shipped with the program, and \btn{Online Help} opens the wiki root. The
bug tracker --- {\scriptsize\file{https://github.com/meerk40t/meerk40t/issues}} --- is also
reachable from the console with \cmd{webhelp issues}, although it is not on the
\menu{Help} menu. The project asks that the tracker be kept for actionable bug
reports and that general questions go to the forums or Discussions.
\end{notebox}

\subsection*{The About window}

\menu{Help>About MeerK40t} opens a window with four tabs. Two of them exist for
bug reports, and you should treat them as part of the tool rather than as
decoration:

\begin{description}[leftmargin=1.4em,style=nextline]
  \item[\tabname{About}] The program version, the project's self-description and
        the credits --- including the people, translators and icon set that made
        the program.
  \item[\tabname{In Memory of David Olsen}] A memorial to David Olsen
        (1982--2024), known as Tatarize, who wrote MeerK40t, with a short history
        of the project.
  \item[\tabname{System-Information}] The exact MeerK40t version, the folder
        holding your configuration file, and a full report of your operating
        system, release, machine and processor. It has \btn{Copy to Clipboard}
        and \btn{Check for Updates} buttons. This is the fastest way to state
        your environment in a bug report.
  \item[\tabname{MeerK40t-Components}] A table of the optional and required
        libraries MeerK40t found --- Python, wxPython, numpy, pillow, the tracing
        engines, serial and USB support, OpenCV, the clipping library and the
        rest --- with a \ui{Version} and a \ui{Status} column and a link to each
        project's home. A component marked \ui{Missing} or \ui{Faulty} explains a
        whole class of ``the feature does not work'' reports, so include this
        table when you file one.
\end{description}

\subsection*{Updates}

Whether MeerK40t looks for new versions is a preference, on the
\tabname{Options} page of the Preferences window, under ``Check for updates on
startup'': \ui{Action} chooses between checking not at all, for major releases,
or for major and beta releases, and \ui{Frequency} chooses between every startup,
once a day or once a week. \menu{Help>Check for Updates} performs the same check
on demand. The answers MeerK40t can give are that a newer release is available,
that you already have the latest version for your system, or that no release
information could be retrieved from GitHub. The same check is available in the
console as \cmd{check_for_updates}, with \cmd{--beta} to include prereleases.

\section*{Reporting a problem so it can be fixed}

A bug report that can be acted on contains five things: what you ran, what you
expected, what happened, the text of the error, and the smallest example that
shows it. Work down this list and you will have all five.

\begin{enumerate}[label=\arabic*.,leftmargin=1.6em]
  \item \textbf{Read the console.} Open the \panel{Console} pane (the ribbon's
        \btn{Console} button, or \menu{Panes>Console}) and scroll back through it
        with the error visible.
        MeerK40t reports failures, warnings about unburnable or unassigned
        elements, and the planning stages here. Copy the text rather than
        paraphrasing it --- a Python traceback's last line and a driver's error
        code are the parts that identify the fault.
  \item \textbf{Note the exact version.} The title bar shows the program name,
        its version and the active device; the About window's
        \tabname{About} and \tabname{System-Information} tabs show the same
        string; \cmd{meerk40t --version} prints it on its own. A build from a git
        checkout appends \ui{git} and the branch, a source tree appends
        \ui{src}, and an installed package appends \ui{pkg} --- that suffix
        matters, because it tells maintainers which artefact you are running.
  \item \textbf{Say which machine and which profile.} Name the device entry from
        the \panel{Devices} pane (the highlighted row is the active one), the
        driver behind it, and how it is connected --- USB, serial, network or
        not connected at all. A fault that only happens on one driver is a
        different report from one that happens with no device at all.
  \item \textbf{Include the diagnostics.} Use \btn{Copy to Clipboard} on the
        \tabname{System-Information} tab, and again on the
        \tabname{MeerK40t-Components} tab, and paste both into the report. If
        the problem is about speed, \cmd{-f}/\cmd{--profiler} records a Python
        profile of a session to a file, which is the useful evidence there.
  \item \textbf{Attach the smallest file that reproduces it.} Delete everything
        from a copy of the project that is not needed to trigger the fault, save
        it, and check that the reduced file still misbehaves. An imported
        original plus your saved project is often enough; a file with one
        rectangle in it is far more likely to be investigated than a
        20\unit{MB} artwork.
  \item \textbf{Describe the steps.} ``Do this, then this, then that'' ---
        including the settings involved (\setting{dpi},
        \setting{bedwidth}, the operation's speed and power) and whether the
        problem repeats on a second attempt. State what you expected to happen
        as well as what did.
  \item \textbf{Search before filing.} Look through the issue tracker and
        Discussions for the same symptom; adding your version and diagnostics to
        an existing report is more useful than opening a near-duplicate.
  \item \textbf{Match your expectations to the project's state.} MeerK40t is in
        a low-maintenance phase: issues and pull requests may wait a long time
        for an answer, and silence usually means lack of time rather than lack of
        interest. Weekly builds are automated snapshots that will contain
        regressions --- if you are chasing a bug in one, say so, and reproduce it
        on the latest stable release before reporting it as a stable-release bug.
\end{enumerate}

\section*{Project status and driver support}

MeerK40t is currently in a low-maintenance phase: the original maintainers have
limited bandwidth for active development, so issues and pull requests may not
receive immediate feedback. Weekly builds are automated experimental snapshots
that will contain bugs and regressions --- if you need a reliable machine, stay
on the latest stable release --- and the fastest route to a fix for something you
find is usually a pull request or a clear, reproducible report.

Hardware support is tracked separately from program features. The table below
reproduces the project's own status list; treat the notes column as the honest
summary of what you are buying into.

\begin{center}\small
\begin{tabular}{@{}p{2.0cm}p{2.6cm}p{2.9cm}p{6.1cm}@{}}
\toprule
\textbf{Device type} & \textbf{Use status} & \textbf{Development status} &
\textbf{Notes} \\
\midrule
Lihuiyu & Minor issues & Under revision & Users report raster Y-axis
registration issues. \\
GRBL & Okay & Stable & No issues known. \\
Moshi & Okay & Stable & No issues known. \\
Newly & Okay & Stable & No issues known. \\
Ruida & Okay & Under revision & Fit for testing (rather than production). \\
Balor & Okay & Stable & No issues known. \\
\bottomrule
\end{tabular}
\end{center}

\noindent Status as published by the project on 17/12/2025. In the project's own
legend, \emph{Okay} means generally reliable for production use, \emph{Minor
issues} means usable but liable to occasional problems, \emph{Stable} means the
core functionality is mature and well tested, and \emph{Under revision} means
active development with changes expected. The driver-specific settings that go
with each of these families are collected in \ref{app:drivers}.

\section*{Final safety checklist}

Tick all of these before pressing \btn{Start}. They take a minute and they cover
the failures that ruin work, burn material or start fires --- in roughly that
order of likelihood.

\begin{center}\small
\begin{tabular}{@{}p{0.5cm}p{13.9cm}@{}}
\toprule
\textbf{} & \textbf{Before the beam comes on} \\
\midrule
$\square$ & Enclosure closed, lid shut, every interlock intact and working ---
no bypassed switches, no tape over sensors. \\
$\square$ & Air assist on, if your machine has it; if it is off, you have
decided that deliberately. \\
$\square$ & Focus checked for this material and this thickness, with the
machine's own gauge or jig. \\
$\square$ & Workpiece secured: clamped, taped or magnetised, and clear of the
head's travel over the whole job. \\
$\square$ & Material identified with certainty. If you do not know what it is,
do not cut it --- PVC, polyurethane, vinyl, ABS and unknown composites release
poisonous fumes. \\
$\square$ & Settings proven on scrap of the same material and thickness, not
assumed from a forum post or a previous job. \\
$\square$ & Bed size and origin correct for this machine, and the job inside the
work area --- the job is not going to run off the bed or into a clamp. \\
$\square$ & Material thickness and kerf accounted for, so that parts will fit
after the cut rather than before it. \\
$\square$ & Job estimated in Simulation or in the \panel{Spooler} estimate
column, and the time and material cost accepted. \\
$\square$ & Nothing flammable inside the enclosure that is not part of the job:
no loose offcuts, no paper, no cleaning rags, no solvent. \\
$\square$ & Extinguisher within reach and suitable for the material, and you
know where the pause and stop controls are --- on the machine and in the
\panel{Laser-Control} pane. \\
$\square$ & You intend to stay with the machine for the whole job. It is not
left unattended, ever. \\
\bottomrule
\end{tabular}
\end{center}
